% !PRINTED 606
\weberpage{606}

The question now concerns the limiting value which these expressions approach when $s=1$.

As regards the left-hand side, it follows from (17) that
\weq{21}{\left\{\begin{aligned}
&\mathrm{Lim}\,(s-1)\,Q_1 = hg \text{ finite and different from zero,}\\
&\mathrm{Lim}\,Q_2\,Q_3\ldots Q_h \text{ finite (perhaps zero)}
\end{aligned}\right.}
is.
To assess the right-hand side we observe that the prime ideals $\mathfrak{p}$ are of the first or of the second degree, and that accordingly $N(\mathfrak{p})$ is equal to a natural prime or to the square of a natural prime. Further, we rely on the well-known theorem of analysis according to which an infinite product
\[
(1-q_1)\,(1-q_2)\,(1-q_3)\ldots
\]
converges unconditionally and has a limiting value different from zero, if
\[
q_1+q_2+q_3+\cdots
\]
is an unconditionally convergent series, and none of the $q_i=1$.

We accordingly divide the prime ideals $\mathfrak{p}$ into two kinds, $\mathfrak{p}_1, \mathfrak{p}_2$, in such a way that the $\mathfrak{p}_1$ are to be of the first degree and to belong to the exponent $\varphi=1$ (hence $\varepsilon=h$). The $\mathfrak{p}_2$ are to be either of the second degree, or to belong to a higher exponent $\varphi$. \gesperrt{Thus the $\mathfrak{p}_1$ are the prime ideals of the first degree belonging to the principal class $\overline{A}_1$.}

Then
\[
N(\mathfrak{p}_1)^{\varphi} = N(\mathfrak{p}_1) = p_1
\]
is a natural prime, and
\[
N(\mathfrak{p}_2)^{\varphi} = p_2^{k}
\]
is at least the 2nd (at most the $2h$th) power of a prime. By the analytic theorem just mentioned, the product
\[
\Pi\bigl[1-N(\mathfrak{p}_2)^{-s\varphi}\bigr]^{\varepsilon}
\]
is then finite and different from zero for $s=1$, and from (20), (21) there follows the theorem:

6. \gesperrt{The limiting value}
\[
\mathop{\mathrm{Lim}}\limits_{s=1} \; \Pi\,\dfrac{s-1}{\Pi\bigl[1-N(\mathfrak{p}_1)^{-s}\bigr]^{h}}
\]
\gesperrt{is finite when $\mathfrak{p}_1$ runs through the prime ideals of the first degree of the principal class $\overline{A}_1$.}

A particular interest attaches to the question whether this limiting value can vanish.
